Modern deep-learning research depends on shared GPU clusters in which
training jobs compete for a~small number of~expensive accelerators.
Operators are expected to~place each job on~the~hardware that meets its
deadline at~the~lowest cost, despite per-GPU prices that vary by~more
than an~order of~magnitude between generations.  Manual scheduling
quickly becomes the~bottleneck, and~off-the-shelf cluster managers
(Kubernetes, plain SLURM) treat GPU bundles as~interchangeable and~have
no~notion of~per-job energy budgets or~deadline-aware preemption.

The ANDREAS project~\cite{andreas} at~Politecnico di~Milano demonstrated
that a~cost-aware optimiser, fed with~profile measurements of~each job
type on~each hardware bundle, can~plan placements that beat the~typical
operator's intuition.  ANDREAS ran on~the~ARMIDA cluster~--- uniform
Tesla~V100 nodes~--- dispatching containers to~them through SLURM, and~its
implementation is~not~available to~reuse.  The~two-node Polimi cluster
targeted here is~instead \emph{heterogeneous}: an~A40 node alongside
a~legacy QuadroP600 node on~different CUDA and~torch generations.  A~job
preempted on~one node must therefore resume across torch versions when
it~migrates to~the~other, and~the~scheduler must consume a~standalone
optimiser (GPUspb~\cite{filippini2023pathrelinking,filippini2024stochastic})
as~an~HTTP service rather than embed one.

This report describes Intelligent Job~Management (IJM), an~independent
reimplementation of~the~ANDREAS approach as~a~distributed control plane.
The~control plane is~a~single FastAPI service that hosts the~profiling scheduler, the
optimiser client, and~the~per-node slot semaphores.  Each GPU node runs
a~lightweight worker that owns the~local Docker daemon.  The~two are
glued together by~plain HTTP over SSH~tunnels, so the~same code runs on
a~laptop talking to~a~remote cluster and~on~a~single-host development
machine.  IJM integrates the~GPUspb optimiser as~the~external decision
maker for~standard-run placement.  Before any instance is~placed as~a
standard run, the~scheduler first executes a~small number of~profile
runs for~it on~previously-unmeasured (node, GPU-count) cells, so the
optimiser is~never asked to~score an~unmeasured bundle.  Migration between
the~A40 and~legacy P600 nodes mid-training is~what makes such placement
worthwhile, and~it~is~not~free: a~PyTorch checkpoint is~normally locked to
the~version that~wrote it, and~the~two nodes can~only exchange one through
a~shared filesystem.  IJM closes both gaps~--- a~deliberately
version-tolerant checkpoint format and~an~rclone-mounted shared store~---
and~asks of~each training job only that~it~honour a~small
\emph{contract}; the~solution and~its reasoning are~detailed in
\S\ref{sec:checkpoints} and~the~contract in~Chapter~\ref{ch:config}.

We validate IJM with~two end-to-end scenarios on~the~two-node cluster,
each exercising the~full job lifecycle and~checking that the~system
faithfully realises the~optimiser's decisions.  Scenario~1
(under-subscribed) submits five instances of~\texttt{cnn\_big}, a
compute-heavy synthetic CNN: four loose-deadline ``patient'' jobs and~one
tight-deadline ``urgent'' job.  It~drives the~profile sweep, single- and
dual-GPU placement, and~an~urgent job that~migrates across nodes~--- off
the~legacy P600 onto~the~A40~--- resuming its checkpoint across torch
versions.  Scenario~2 fills the~cluster and~mixes types (\texttt{cnn\_big}
pins on~A40, slack \texttt{lstm-small} patients~--- a~small LSTM~--- on
P600), adding deadline-driven preemption of~the~cheapest slack work and~a
second cross-node migration.  Together they confirm that profiling,
placement, preemption, and~cross-version resume behave correctly and~stay
consistent with~the~optimiser's plan.  (As a~secondary observation, the
runs also expose a~\emph{horizon-myopia} in~the~GPUspb \texttt{method=RG}
proxy~--- an~urgent job parked on~slow-but-free hardware~--- which
Chapter~\ref{ch:e2e} characterises but~which does not~affect IJM's own
correctness.)

Chapter~\ref{ch:architecture} describes the~deployment topology, the
scheduling pipeline, the~state machine, the~persistent state, and~the
invariants that protect the~control plane from concurrency hazards.
Chapter~\ref{ch:config} documents the~cluster and~job configuration
files, the~submission API, and~the~runtime knobs.
Chapter~\ref{ch:e2e} presents the~two scenarios and~the~cost model used
by~the~optimiser.
Chapter~\ref{ch:conclusion} summarises the~contributions, the~limitations,
and~the~follow-up work that has concrete entry points in~the~codebase.
